\section{Background and Requirements}
\label{sec:background}

\paragraph{Governance by architecture.}
SARC~\cite{besanson2026sarc} treats constraints as first-class
specification objects alongside state, action space, and reward: a
constraint declares its source, class (\emph{hard} / \emph{soft} /
\emph{escalation}), predicate, verification point, and response, and
compiles into four enforcement points in the agent loop — a Pre-Action
Gate, an Action-Time Monitor, a Post-Action Auditor, and an Escalation
Router — bound by invariants whose joint effect is a \emph{decidable
audit}: given a specification \Sig with constraint set $C$ and a
trace $T$, a checker decides $T \entails \Sig$ — read: the trace
satisfies the specification — in $O(|T|\cdot|C|)$, trace records
times constraints, without access to the model or its prompts. SARC stops at the loop: its reference artifact keeps \Sig
in a file beside the system and audits an exported trace with an
external checker. The knowledge the agents act on sits in an
ungoverned store underneath. Our position is that the store is the
right compilation target for exactly the machinery SARC specifies.

\paragraph{Bitemporal data.}
A bitemporal store indexes facts by transaction time (when the store
learned something) and valid time (when it holds in the world), in the
lineage of Datomic and XTDB. What conventional bitemporal systems do
not do is extend the treatment beyond data: trust annotations, policy
decisions, and the validation rules themselves remain latest-only, so
the historical questions governance actually asks — what was believed,
what was trusted, what was \emph{required} at $T$ — outrun the store's
memory.

\paragraph{Named graphs and trust.}
RDF datasets partition triples into named graphs, and the provenance
literature attached trust semantics to them two decades
ago~\cite{carroll2005named}; information-flow lattices are older
still~\cite{denning1976lattice}. SPARQL's dataset construction,
however, is silent union: nothing in the standard machinery refuses a
composition, degrades it, or even remarks that one member graph carries
no declaration at all.

\paragraph{Requirements.}
A store serving multi-writer agent ingestion in a governed setting
needs, at minimum: write-time constraint evaluation against the state
the write would create; a permanent, attributable record of every
gate decision, including refusals; authority that attaches to
partitions and only narrows under delegation; composition that cannot
launder trust; an audit decidable from the store's own contents; and
reproducibility of past decisions under the rules in force at the
time. \S\ref{sec:principles} states these as six principles;
\S\ref{sec:datamodel}--\S\ref{sec:governance} show the mechanisms that
discharge them.
